\begin{table}[htbp]
\centering
\caption{Placebo test: the same regression run on every last price digit}
\label{tab:placebo}
\small
\begin{threeparttable}
\begin{tabular}{@{}lrrr@{}}
\toprule
Last digit & Coefficient & SE & $t$ \\
\midrule
0 & 0.0192*** & (0.0075) & 2.58 \\
1 & -0.0128 & (0.0092) & -1.40 \\
2 & 0.0057 & (0.0103) & 0.55 \\
3 & 0.0043 & (0.0085) & 0.50 \\
4 & -0.0023 & (0.0099) & -0.24 \\
5 & -0.0017 & (0.0092) & -0.19 \\
6 & -0.0008 & (0.0097) & -0.09 \\
7 & -0.0088 & (0.0094) & -0.94 \\
8 & -0.0120 & (0.0091) & -1.32 \\
9 & -0.0089 & (0.0098) & -0.91 \\
\bottomrule
\end{tabular}
\begin{tablenotes}[flushleft]\footnotesize
\item Outcome is the log effective spread in the dynamic specification of Table~\ref{tab:rq2}, with the round-price share replaced by the share of buy-initiated large-trade volume at each last digit in turn. If the round-price result reflects stale limit orders at psychologically salient prices, digits other than zero should carry no comparable information. Digit five is marked as semi-focal because half-way points attract some of the same behaviour.
\end{tablenotes}
\end{threeparttable}
\end{table}
